\scanpage{68}
\[
(1)\qquad\Pi_1(X)\unsure{\relbar\joinrel\relbar|\longrightarrow}\widehat\Pi_1(X)
\]
qui induit une bijection sur l'ens des objets \inserted{ou ``sommets''} -- s'identifiant à $X=X(\mathbb C)$ -- et qui \deleted{déf} pour $\forall x\in X=X(\mathbb C)$, donne l'hom de groupes
\[
(2)\qquad
\begin{array}{ccc}
\operatorname{Aut}_{\Pi_1(X)}(x)&\longrightarrow&\operatorname{Aut}_{\widehat\Pi_1(X)}(x)\\
\Vert\ \text{déf}&&\Vert\ \text{déf}\\
\pi_1(X,x)&&\widehat\pi_1(X_{\mathrm{sch}},x)
\end{array}
\]
permettant d'identifier le $\pi_1$ algébrique au compactifié profini du $\pi_1$ ``topologique'' ou ``transcendant'' : \unsure{b}
\[
(3)\qquad\pi_1(X_{\mathrm{sch}},x)\simeq\widehat\pi_1(X,x).
\]
De $=$, pour $x,y\in X$, l'application
\[
\operatorname{Hom}_{\Pi_1(X)}(x,y)\longrightarrow\operatorname{Hom}_{\widehat\Pi_1(X)\deleted{\unsure{y}}}(x,y)
\]
\margin{On identifiera un élément de $\operatorname{Hom}_{\Pi_1}(x,y)$ avec son image dans $\operatorname{Hom}_{\widehat\Pi_1}(x,y)$, sans se préoccuper si l'application est injective ou non.}
permet d'identifier le deuxième ens, en tant que bi-torseur sur $(\widehat\pi_1(X,y),\widehat\pi_1(X,x))$ au bitorseur déduit \deleted{Hom} du premier ens par extension des groupes d'opérateurs par
\[
\pi_1(X,y)\unsure{\relbar\joinrel\relbar|\longrightarrow}\widehat\pi_1(X,y),\qquad\pi_1(X,x)\unsure{\relbar\joinrel\relbar|\longrightarrow}\widehat\pi_1(X,x).
\]
(ce couple d'homomorphismes satisfaisant \inserted{bien} à la condition requise pour que cette extension ait un sens\ldots).
\margin{à savoir (en termes de $(G,H),(G',H')$ et $G\to G',H\to H'$ : il existe un isom ext. $G'\xrightarrow{\sim}H'$ rendant commutatif le diagramme \deleted{\unsure{des bitorseurs}}
\diagram{\begin{tikzcd}[ampersand replacement=\&,column sep=large,row sep=large]
G\arrow[r,"\sim", "\substack{\text{défini par le}\\\text{bitorseur }(G,H)}"']\arrow[d]\&H\arrow[d]\\
G'\arrow[r,"\sim"]\&H'
\end{tikzcd}}}
Soit
\[
(5)\qquad\unsure{\mathbf{\Gamma}|\Gamma}_{\deleted{K}}=\unsure{\mathbf{\Gamma}|\Gamma}_K=\operatorname{Aut}_{K\text{-cor}}(\mathbb C),
\]
\margin{67}
